\documentclass[border=7bp]{standalone}
% Shared preamble for the report's editable vector figures.
\usepackage{fontspec}
\usepackage{xeCJK}
\setmainfont{Calibri}
\setCJKmainfont{SimSun}[AutoFakeBold=1.5]
\usepackage{amsmath}
\usepackage{tikz}
\usetikzlibrary{arrows.meta,positioning,calc,shapes.geometric,fit}
\definecolor{Ink}{HTML}{25313D}
\definecolor{Blue}{HTML}{315E80}
\definecolor{PaleBlue}{HTML}{E7EFF5}
\definecolor{Muted}{HTML}{626B74}
\definecolor{PaleGray}{HTML}{F0F1F3}
\tikzset{
  >=Stealth,
  every node/.style={font=\fontsize{10bp}{13bp}\selectfont,text=Ink},
  flow/.style={->,draw=Ink,line width=.65pt},
  frame/.style={rectangle,draw=Ink,line width=.65pt,align=center,
    minimum height=.85cm,text width=5cm,inner sep=5pt},
  emphasis/.style={frame,draw=Blue,fill=PaleBlue},
  supporting/.style={frame,draw=Muted,fill=PaleGray},
  annotation/.style={font=\fontsize{9bp}{12bp}\selectfont,text=Muted,align=center},
}

\begin{document}
\begin{tikzpicture}
\path (-7.3,0) rectangle (7.3,0);

\node[emphasis,text width=11cm] (input) at (0,0) {固定输入、算法、启发函数、堆上限和缓存条件};
\node[frame,text width=11cm] (start) at (0,-1.3) {父进程开始计时，启动独立工作 JVM};
\node[frame,text width=11cm] (prepare) at (0,-2.6) {读取输入并准备启发函数\\冷 PDB：建表；热 PDB：校验与加载；其他启发：无建表};
\node[emphasis,text width=11cm] (search) at (0,-4.1) {搜索计时：调用搜索器 → 返回结果\\包含启发查询和同步快照记录};
\node[frame,text width=4cm] (ok) at (-2.7,-5.8) {SOLVED\\恢复路径、逐步回放、导出};
\node[supporting,text width=4cm] (fail) at (2.7,-5.8) {TIMEOUT / OOM / FAIL\\保留失败记录与可用快照};
\node[frame,text width=11cm] (end) at (0,-7.4) {工作进程退出或终止回收，父端计时结束};
\node[supporting,text width=11cm] (stats) at (0,-8.8) {分别统计成功次数、搜索与父端中位数\\失败不填作成功时间，不混入成功中位数};
\foreach \a/\b in {input/start,start/prepare,prepare/search,end/stats}{\draw[flow](\a)--(\b);}
\draw[flow](search.south)-|(ok.north); \draw[flow](search.south)-|(fail.north);
\coordinate (join) at (0,-6.7);
\draw[draw=Ink,line width=.65pt](ok.south)|-(join);
\draw[draw=Ink,line width=.65pt](fail.south)|-(join);\draw[flow](join)--(end.north);

\end{tikzpicture}
\end{document}
